\pdfoutput=1
\documentclass[11pt]{article}

\usepackage[margin=1.05in]{geometry}
\usepackage[utf8]{inputenc}
% Latin Modern with T1 when available (arXiv has it); plain CM otherwise
\IfFileExists{lmodern.sty}{\usepackage[T1]{fontenc}\usepackage{lmodern}}{}
\usepackage{amsmath,amssymb,amsthm}
\usepackage{booktabs}
\usepackage{microtype}
\usepackage{enumitem}
\usepackage[hidelinks]{hyperref}

\renewcommand{\abstractname}{R\'esum\'e}
\renewcommand{\refname}{R\'ef\'erences}
\renewcommand{\proofname}{Preuve}

\newtheorem{theorem}{Th\'eor\`eme}
\newtheorem{corollary}{Corollaire}
\newtheorem{proposition}{Proposition}
\newtheorem{lemma}{Lemme}
\theoremstyle{definition}
\newtheorem{definitionx}{D\'efinition}

% Shared macro definitions for paper.tex and paper.fr.tex.
% One source of truth: a symbol renamed here is renamed in both papers.
\newcommand{\Fp}{\mathbb{F}_p}
\newcommand{\negl}{\mathsf{negl}}
\newcommand{\LeafHash}{\mathsf{LeafHash}}
\newcommand{\Compress}{\mathsf{Compress}}
\newcommand{\PRF}{\mathsf{PRF}}
\newcommand{\CtxHash}{\mathsf{CtxHash}}
\newcommand{\Trunc}{\mathsf{Trunc}_4}
\newcommand{\MerkleFold}{\mathsf{MerkleFold}}
\newcommand{\RO}{\mathsf{RO}}
\newcommand{\Enc}{\mathsf{Enc}}
% Poseidon2 state width is t, so the Fiat-Shamir challenge vector is
% chi; the per-threshold accumulators are calligraphic, so they do not
% collide with the transcript body T.
\newcommand{\chall}{\chi}
\newcommand{\tree}{\mathcal{T}}
\newcommand{\adv}{\mathcal{A}}
% := without pulling in mathtools; arXiv has it, but one fewer package
% is one fewer thing that can differ between their TeX Live and yours.
\providecommand{\coloneqq}{\mathrel{\vcenter{\hbox{$:$}}}\!=}


\title{Appartenance anonyme à un ensemble, post-quantique,\\
pour une vérification d'âge en ligne sans surveillance}
\author{Alexandre `kidev' Poumaroux}
\date{}

\begin{document}
\maketitle

\begin{abstract}
Nous présentons un protocole qui permet aux sites web de vérifier qu'un visiteur a atteint
la majorité tout en rendant cryptographiquement impossible, pour le site vérifieur comme
pour le gouvernement émetteur, de savoir qui est ce visiteur ou quels sites un citoyen
donné fréquente. La construction centrale est l'appartenance anonyme à un ensemble : le
gouvernement publie une unique racine de Merkle signée engageant l'ensemble des
engagements de hachage des adultes vérifiés, et le portefeuille open source de chaque
citoyen prouve à divulgation nulle de connaissance qu'il connaît la préimage d'une feuille
sous cette racine, ainsi que la dérivation correcte d'une clé pseudonyme propre à chaque
site et à chaque époque. La vérification est effectuée entièrement hors ligne par le site,
contre la racine publiée ; le gouvernement est structurellement absent de chaque flux
d'usage, de sorte qu'il n'existe aucun trafic de vérification, aucun horodatage et aucun
compteur qui puisse être journalisé, corrélé ou réquisitionné. Toutes les primitives sont
fondées sur le hachage ou sur les réseaux euclidiens : Poseidon2 sur le corps de
Goldilocks pour les engagements, la compression de Merkle et la dérivation des pseudonymes
; le système de preuve MPC-in-the-head ZKBoo compilé par Fiat-Shamir pour l'argument à
divulgation nulle ; et ML-DSA-65 (FIPS 204) pour l'authentification de la racine. La
conception résiste donc aux adversaires quantiques. Nous donnons des définitions par jeux
et des preuves complètes pour l'anonymat, l'infalsifiabilité et la résistance au rejeu, un
traitement honnête de la seule propriété qu'aucun système d'accréditations ne peut imposer
(le prêt volontaire), et les performances mesurées d'une implémentation de référence
complète : à 128 bits de solidité sur un arbre d'une capacité de 4,29 milliards de
feuilles, la génération de preuve prend 4,7\,s et la vérification 3,0\,s sur un seul c\oe
ur d'un AMD Ryzen~9 7950X en Python pur, pour une preuve de 21,8\,Mo, tandis que le portage Rust du
même moteur qui accompagne cet article prouve en 0,33\,s et vérifie en 0,18\,s sur un c\oe
ur, quinze fois plus vite que CPython mesuré sur cette même machine. Nous analysons la
trajectoire vers des preuves sous les 200\,Ko et une vérification en dizaines de
millisecondes avec des moteurs de preuve de production.
\end{abstract}

\section{Introduction}

Les obligations de vérification d'âge en ligne se multiplient, pour les réseaux sociaux,
les magasins d'applications et un nombre croissant de services soumis à un contrôle d'âge,
et chaque mécanisme de conformité largement déployé impose un choix fatal : soit le site
apprend qui vous êtes (envoi de documents, scan du visage), soit un tiers apprend où vous
allez (services de vérification centralisés, identité fédérée). Les deux créent des bases
de données dont l'existence même constitue le préjudice : un registre permanent de qui
visite quoi en ligne est une infrastructure de surveillance générale, à une fuite, une
réquisition ou un changement de politique près d'être retournée contre n'importe quel
citoyen ou groupe. Le résultat est prévisible : les adultes contournent les contrôles, les
mineurs empruntent les mêmes voies de contournement, et l'objectif affiché de protection
de l'enfance échoue tandis que le risque de surveillance demeure.

Nous prenons pour contrainte de conception la version la plus forte possible de l'exigence
de vie privée, posée comme règle inviolable : le système doit permettre à un site de
confirmer que ce visiteur est un adulte certifié par le gouvernement, sans donner ni au
site, ni au gouvernement, ni aux deux réunis, aucune capacité, jamais, de relier une
visite à un citoyen. La garantie doit être mathématique et non institutionnelle : un
citoyen qui suppose le gouvernement activement malveillant doit rester en sécurité en
utilisant le système.

Formellement, le système doit satisfaire :
\begin{description}[leftmargin=2.6em,style=nextline]
\item[R1 (Solidité).] Un site n'accepte que les visiteurs détenant une accréditation
délivrée par le gouvernement à un adulte vérifié.
\item[R2 (Aveuglement de l'émetteur).] Le gouvernement ne peut déterminer quel citoyen se
trouve derrière un événement de vérification, ni même qu'un citoyen donné a utilisé le
système après son inscription.
\item[R3 (Aveuglement du vérifieur).] Le site apprend exactement un bit (majeur : oui)
plus, en option, un pseudonyme stable uniquement au sein de ce site et d'une fenêtre
temporelle bornée ; les pseudonymes entre sites sont non chaînables, même pour des sites
en collusion.
\item[R4 (Résistance à la collusion).] R2 et R3 tiennent même si le gouvernement et un
ensemble quelconque de sites mettent toutes leurs données en commun.
\item[R5 (Sécurité post-quantique).] Toutes les propriétés tiennent face à des adversaires
quantiques.
\end{description}

Deux décisions structurelles font l'essentiel du travail. Premièrement, le gouvernement ne
signe jamais aucune valeur propre à un citoyen : sa production cryptographique entière est
une signature sur un unique accumulateur public, si bien qu'aucune signature ne peut
jamais servir de marqueur de traçage. Deuxièmement, la vérification n'exige que des
paramètres publics, elle est donc effectuée localement par le site ; le gouvernement n'a
aucun point de terminaison de vérification, et des messages jamais envoyés ne peuvent être
ni journalisés, ni corrélés, ni réquisitionnés.

Nos contributions :
\begin{enumerate}[leftmargin=2em]
\item Un protocole fondé sur l'appartenance anonyme à un ensemble qui satisfait R1 à R5,
dans lequel l'empreinte en ligne totale du gouvernement est une seule valeur publique
signée, et où la vérification est un calcul local du site (Section~\ref{sec:protocol}).
\item Des définitions de sécurité par jeux et des preuves complètes
(Section~\ref{sec:security}), avec un décompte exact de ce que chaque partie et chaque
coalition en collusion peut calculer, et un énoncé explicite des propriétés que le système
ne fournit pas.
\item Une implémentation de référence complète aux paramètres de production :
Poseidon2-Goldilocks avec les constantes officielles de ses concepteurs, des empreintes de
256 bits, un accumulateur de Merkle creux de capacité $2^{32}$, une solidité Fiat-Shamir
de $2^{-128}$, une signature de racine ML-DSA-65, un format de portefeuille chiffré et un
format binaire de transport ; un portage Rust natif du prouveur et du vérifieur,
interopérable avec elle au niveau du format de transport ; ainsi que des mesures de
performance et une suite de tests
couvrant les attaques analysées en Section~\ref{sec:security} (Sections~\ref{sec:impl},
\ref{sec:perf}).
\end{enumerate}

\section{Travaux connexes}

\paragraph{Accréditations anonymes.} Les signatures aveugles de Chaum \cite{chaum} ont
ouvert l'étude des accréditations que l'émetteur ne peut tracer. Brands \cite{brands} et
Camenisch-Lysyanskaya \cite{cl} ont construit des systèmes complets d'accréditations
anonymes ; les déploiements modernes (Idemix, BBS+ dans les profils de verifiable
credentials \cite{w3cvc}) descendent des seconds. Toutes les instanciations pratiques
reposent sur RSA, le logarithme discret ou les couplages et tombent face à l'algorithme de
Shor \cite{shor} ; les accréditations anonymes et signatures aveugles post-quantiques
restent des prototypes de recherche aux clés volumineuses, aux signatures volumineuses, à
la sécurité concurrente restreinte ou aux hypothèses non standard \cite{aksy,delpino}.

\paragraph{Pseudonymie par appartenance à un ensemble.} Notre construction descend
architecturalement du signalement anonyme à la Semaphore \cite{semaphore} et de la
conception des nullifiers de Zerocash \cite{zerocash} : l'identité est un secret dont
l'engagement vit dans un ensemble de Merkle public, l'appartenance est prouvée à
divulgation nulle, et des sorties de PRF déterministes fournissent des pseudonymes à
portée limitée (nullifiers externes). Nous adaptons ce schéma à la vérification d'âge,
remplaçons ses systèmes de preuve et ses hachages fondés sur le logarithme discret par des
composants post-quantiques, et ajoutons une liaison de transcription côté vérifieur pour
la résistance au rejeu sans aucun serveur de coordination.

\paragraph{Jetons anonymes à débit limité.} Privacy Pass \cite{privacypass} et ses
variantes à débit limité permettent de dépenser des jetons non chaînables ; ils traitent
un problème différent (non-chaînabilité par requête avec un émetteur en ligne) et leurs
fondations VOPRF sont pré-quantiques.

\paragraph{Systèmes de vérification d'âge.} Les prestataires commerciaux (envoi de
documents, estimation faciale) révèlent identité ou biométrie à une partie qu'il faut
croire sur parole. Le modèle W3C Verifiable Credentials \cite{w3cvc} permet la divulgation
sélective, mais ses profils de présentation non chaînable reposent sur BBS+ (couplages,
donc pré-quantique), et ses flux de présentation impliquant l'émetteur laissent fuir des
métadonnées d'usage. Les travaux d'attestation d'âge du portefeuille d'identité numérique
européen partagent notre objectif pour R3 mais pas notre modèle de menace pour R2/R4 : ils
supposent que l'émetteur s'abstient de corréler.

\paragraph{Systèmes de preuve.} Les preuves à divulgation nulle de connaissance remontent
à Goldwasser, Micali et Rackoff \cite{gmr}. ZKBoo \cite{zkboo} et ses successeurs
(ZKB++/Picnic \cite{picnic}, KKW \cite{kkw}) construisent des arguments à divulgation
nulle dans le paradigme MPC-in-the-head d'Ishai et al.\ \cite{ikos}, avec des primitives
exclusivement symétriques : aucune mise en place de confiance, sécurité post-quantique
plausible. Les STARK \cite{stark} atteignent le même profil de confiance avec des preuves
succinctes au prix d'une complexité d'implémentation. Notre implémentation de référence
utilise ZKBoo pour son auditabilité ; la Section~\ref{sec:deploypath} quantifie la
trajectoire de déploiement STARK/VOLE-in-the-head.

\paragraph{Standards post-quantiques.} Le NIST a standardisé ML-DSA (Dilithium) sous FIPS
204 \cite{fips204,dilithium} et SHA-3 sous FIPS 202 \cite{fips202} ; Poseidon2
\cite{poseidon2} est la permutation adaptée à l'arithmétisation de référence pour les
systèmes à divulgation nulle. L'algorithme de Grover \cite{grover} divise par deux les
niveaux de sécurité symétriques, ce que nos choix de paramètres prennent en compte.

\section{Système et modèle de menace}
\label{sec:model}

\paragraph{Participants.} Une Autorité Gouvernementale (AG) qui vérifie l'âge par les
mécanismes existants d'état civil ; des Citoyens exécutant un Portefeuille open source sur
leurs propres appareils ; des Sites web qui doivent conditionner leur contenu à la
majorité.

\paragraph{Hypothèses de confiance.} (i) L'AG vérifie correctement l'âge à l'inscription
et n'inscrit que des adultes ; (ii) les primitives cryptographiques satisfont leurs
notions de sécurité standard (résistance aux collisions et sécurité PRF des fonctions
fondées sur Poseidon2 de la Section~\ref{sec:prims}, EUF-CMA de ML-DSA-65, comportement
d'oracle aléatoire de SHA3/SHAKE \cite{bellrog} pour les engagements et Fiat-Shamir) ;
(iii) l'appareil du citoyen n'est pas compromis. Nous ne supposons explicitement
\emph{pas} que l'AG s'abstient de corréler, de s'entendre ou de contraindre ; que les
sites se comportent bien ; ni que les citoyens gardent leurs accréditations pour eux.

\paragraph{Adversaires.} Une AG malveillante peut enregistrer tout ce qu'elle voit,
s'entendre avec des sites et tenter des attaques par vue divisée (split-view) sur les
données publiées. Des sites malveillants peuvent mettre leurs données en commun entre eux
et avec l'AG et tenter un chaînage entre sites. Des citoyens malveillants peuvent tenter
de forger une appartenance, de rejouer des jetons ou de prêter leurs accréditations. Tous
les adversaires peuvent disposer de capacités quantiques.

\paragraph{Objectifs de sécurité.}
\textbf{G1} Infalsifiabilité : aucune coalition dépourvue d'un secret inscrit ne produit
un jeton accepté, sauf avec probabilité négligeable.
\textbf{G2} Anonymat : pour tout jeton accepté, aucune coalition de l'AG et de sites ne
peut distinguer lequel des citoyens inscrits l'a produit avec un avantage non négligeable
sur le tirage au hasard parmi tous les citoyens inscrits compatibles avec les données
publiques.
\textbf{G3} Non-chaînabilité entre sites : des jetons présentés à des sites différents (ou
à des époques différentes) ne peuvent être reliés à un même citoyen.
\textbf{G4} Résistance au rejeu : un jeton accepté une fois n'est plus accepté, ni sur le
même site ni sur un autre.
\textbf{G5} Tout ce qui précède face à des adversaires quantiques.

\paragraph{Hors périmètre.} Les identifiants de la couche réseau (adresses IP) sont
visibles des sites, comme pour toute requête HTTPS ; les citoyens qui exigent l'anonymat
réseau composent le système avec Tor ou un VPN, avec lesquels il n'interfère pas,
précisément parce qu'aucun point de terminaison gouvernemental n'est contacté. La
compromission de l'appareil et la contrainte légale exercée sur le citoyen lui-même sont
hors périmètre.

\section{Le protocole}
\label{sec:protocol}

\subsection{Primitives et paramètres}
\label{sec:prims}

\paragraph{Corps.} Le premier de Goldilocks $p = 2^{64} - 2^{32} + 1$. Toutes les valeurs
du protocole sont des vecteurs sur $\Fp$ ; les empreintes font 4 éléments (256 bits).

\paragraph{Permutation.} Poseidon2 \cite{poseidon2} avec un état de largeur $t = 8$, une
S-box $x^7$, 8 tours externes et 22 tours internes, avec les constantes de tour et la
matrice interne publiées par ses concepteurs. Dans tout l'article, $P : \Fp^8 \to \Fp^8$
désigne cette permutation et $\Trunc : \Fp^8 \to \Fp^4$ la projection d'un état sur ses
quatre premiers éléments. De $P$ nous dérivons trois fonctions à séparation de domaine, au moyen de trois
constantes de corps \emph{distinctes} $\mathit{tag}_{\mathrm{leaf}},
\mathit{tag}_{\mathrm{prf}}, \mathit{tag}_{\mathrm{empty}} \in \Fp$ (c'est leur distinction
deux à deux qui rend la séparation de domaine porteuse dans le
Théorème~\ref{thm:unforge}) :
\begin{align*}
\LeafHash(s) &= \Trunc\bigl(P(s_0, s_1, s_2, s_3, \mathit{tag}_{\mathrm{leaf}}, 0, 0, 0)\bigr),\\
\Compress(l, r) &= \Trunc\bigl(P(l \,\|\, r) + (l \,\|\, r)\bigr),\\
\PRF(s, c) &= \Trunc\bigl(P(P(s_0, s_1, s_2, s_3, \mathit{tag}_{\mathrm{prf}}, 0, 0, 0) \boxplus c)\bigr),
\end{align*}
soit l'engagement d'inscription du secret $s \in \Fp^4$, la compression de Merkle 2 vers 1
(le mode de compression avec réinjection (feed-forward) recommandé par les auteurs de
Poseidon2) et la dérivation de pseudonyme. Ici $l \,\|\, r \in \Fp^8$ note la
concaténation et $x \boxplus c \coloneqq (x_0 + c_0, \ldots, x_3 + c_3, x_4, \ldots, x_7)$
ajoute le contexte $c$ de 4
éléments dans le débit. Les contextes sont $c = \CtxHash(\mathrm{domaine},
\mathrm{\acute{e}poque})$, 4 éléments dérivés par SHAKE256 avec échantillonnage par rejet
; les déploiements multi-seuils incorporent le seuil d'âge dans la chaîne de domaine
hachée (Section~\ref{sec:thresholds}).
Sous les analyses standard des modes éponge et compression, ces fonctions offrent 128 bits
de résistance aux collisions et 128 bits de résistance post-quantique à la préimage
(capacité et empreinte font 256 bits ; Grover \cite{grover} divise l'exposant par deux).

\paragraph{Accumulateur.} Un arbre de Merkle incrémental creux \cite{merkle} de profondeur
$d = 32$ (capacité $2^{32} \approx 4{,}29$ milliards, suffisante pour toute population
nationale ou fédérée) avec des feuilles vides toutes égales $E =
\Trunc(P(0,0,0,0,\mathit{tag}_{\mathrm{empty}},0,0,0))$ et des empreintes de sous-arbres
vides mises en cache, de sorte que l'insertion et le calcul de chemin coûtent $O(d)$ quel
que soit le remplissage.

\paragraph{Argument à divulgation nulle.} ZKBoo \cite{zkboo}, un système d'argument dans
le paradigme MPC-in-the-head \cite{ikos} réalisant les preuves à divulgation nulle
\cite{gmr}, sur le circuit arithmétique de la Section~\ref{sec:issuance}, avec la
décomposition additive $(2,3)$ sur $\Fp$, des rubans (tapes) de partie développés depuis
des graines de 256 bits par SHAKE256 en mode compteur, des engagements de vue SHA3-256 et
des défis Fiat-Shamir dérivés par SHA3-256 sur la transcription complète, modélisée comme
un oracle aléatoire \cite{bellrog}. Chaque répétition a une erreur de solidité de $2/3$ ;
nous utilisons $\tau = 219$ répétitions, soit une erreur $(2/3)^{219} < 2^{-128}$. Les
parts d'entrée de deux parties sont dérivées des rubans (l'optimisation ZKB++
\cite{picnic}) ; les parts de la troisième partie et les sorties de multiplication de la
vue voisine sont les seules données transmises par porte.

\paragraph{Authentification de la racine.} ML-DSA-65 (FIPS 204 \cite{fips204}), catégorie
de sécurité NIST 3, signe chaque lot de racine publié.

\paragraph{Stockage du portefeuille.} Le secret est chiffré au repos avec AES-256-GCM sous
une clé dérivée par scrypt ($N = 2^{15}$, $r = 8$, $p = 1$).

\subsection{Inscription}

Le portefeuille tire $s \leftarrow \Fp^4$ uniformément (256 bits d'entropie locale qui ne
quittent jamais l'appareil) et calcule $\mathit{leaf} = \LeafHash(s)$. Le citoyen
s'authentifie auprès de l'AG par le mécanisme d'identité civile existant de sa juridiction
(en personne ou par eID) et ne soumet que $\mathit{leaf}$. L'AG vérifie que le citoyen est
un adulte non déjà inscrit, ajoute $\mathit{leaf}$ à l'arbre à l'indice libre suivant et
enregistre $(\mathrm{id\_citoyen}, \mathrm{indice})$. La génération de clé est entièrement
locale : produire $s$ et $\mathit{leaf}$ n'exige aucune connexion réseau, si bien qu'une
application portefeuille peut créer son accréditation complètement hors ligne (même sur un
appareil isolé du réseau) et ne divulguer l'engagement qu'au moment de l'inscription ;
rien dans la clé ne dépend jamais d'un serveur.

L'AG a le droit de retenir cette correspondance pour toujours. C'est sans danger :
$\mathit{leaf}$ et l'indice n'apparaissent dans aucun message ultérieur que l'AG pourrait
observer. L'anonymat ne repose pas sur l'oubli de l'AG ; il repose sur le fait que l'AG ne
voit jamais rien d'autre.

\subsection{Publication de la racine}
\label{sec:rootpub}

Après chaque lot d'inscriptions ou de révocations, l'AG publie un lot de racine (root
bundle)
\[
B = \bigl(\mathit{racine}, \mathit{taille}, \mathit{emis}, d,\
\mathrm{Sig}_{\text{ML-DSA-65}}(\ell \,\|\, \mathit{racine} \,\|\, \mathit{taille} \,\|\,
\mathit{emis} \,\|\, d)\bigr),
\]
où $\mathit{taille}$ est le nombre de feuilles allouées à ce jour, $\mathit{emis}$
l'horodatage de publication, et $\ell$ une étiquette de version fixe (étendue par le seuil d'âge dans le
déploiement multi-seuils de la Section~\ref{sec:thresholds}), publié avec la séquence
complète des feuilles, l'arbre étant une donnée publique. Les
lots sont distribués à la manière d'un journal de transparence : miroirs, diffusion et
audit en ajout seul, de sorte que l'AG ne peut servir des racines différentes à des
utilisateurs différents (une attaque par vue divisée lui permettrait sinon de donner à un
citoyen un arbre personnalisé et de reconnaître les preuves qui en résultent ; la
diffusion réduit cela à une attaque détectable en comparant deux observations de lots).
Les portefeuilles synchronisent l'arbre en bloc, ou par récupération privée d'information,
de sorte qu'obtenir son chemin ne révèle rien sur la feuille qui est la sienne ; les sites
n'ont besoin que du lot.

\subsection{Émission du jeton}
\label{sec:issuance}

Pour visiter le domaine $D$ à l'époque $e = \lfloor \mathrm{maintenant} / 30\
\mathrm{jours} \rfloor$, le portefeuille obtient du site un nonce frais $n$ de 128 bits,
calcule $c = \CtxHash(D, e)$, le pseudonyme $k = \PRF(s, c)$, le chemin de Merkle $(b_0,
\sigma_0, \ldots, b_{d-1}, \sigma_{d-1})$ de sa feuille, où $b_i \in \{0,1\}$ est le
$i$-ième bit de l'indice de la feuille et $\sigma_i \in \Fp^4$ l'empreinte frère au niveau
$i$, et une preuve $\pi$ de l'énoncé
\begin{align*}
S(\mathit{racine}, k, c):\quad \exists\, s, (b_i, \sigma_i)_{i<d}:\quad
& b_i \in \{0,1\},\\
& \MerkleFold\bigl(\LeafHash(s), (b_i, \sigma_i)_{i<d}\bigr) = \mathit{racine},\quad
\PRF(s, c) = k.
\end{align*}
Ici, $\MerkleFold$ est le recalcul itéré de racine défini par $h_0 = x$ et, pour $0 \le i
< d$, $h_{i+1} = \Compress(h_i, \sigma_i)$ si $b_i = 0$ et $h_{i+1} = \Compress(\sigma_i,
h_i)$ si $b_i = 1$, avec $\MerkleFold\bigl(x, (b_i, \sigma_i)_{i<d}\bigr) = h_d$. L'énoncé
est réalisé comme circuit arithmétique : une permutation $\LeafHash$ ; par niveau, un
échange conditionnel à 1 multiplication par élément d'empreinte, une contrainte de
booléanité $b_i^2 - b_i = 0$ exposée comme sortie publique nulle, et une permutation
$\Compress$ ; puis deux permutations pour la PRF. La taille du circuit est $344(d+3) + 5d$
portes de multiplication (12\,200 pour $d = 32$). La transcription Fiat-Shamir hache, en
plus de tous les engagements et parts de sortie, la chaîne de contexte
\[
\mathit{ctx} = \text{``avsm-token/1''} \,\|\, D \,\|\, e \,\|\, \mathit{racine} \,\|\, k
\,\|\, n,
\]
liant la preuve au domaine, à l'époque, à la racine, au pseudonyme revendiqué et au nonce
du site, les champs consécutifs étant joints par un octet séparateur réservé fixe (l'octet
\texttt{|}, absent de tout nom d'hôte). Notons $\Enc(D, e, \mathit{racine}, k, n)$ cette
chaîne d'octets ; l'argument de rejeu de la Section~\ref{sec:security} en exige une
propriété syntaxique.

\begin{lemma}[Encodage du contexte]
\label{lem:enc}
Si le champ de domaine ne contient aucune occurrence de l'octet séparateur, alors $\Enc$
est injective.
\end{lemma}

\begin{proof}
Tous les champs qui suivent le domaine ont une largeur fixe : 8 octets pour $e$, 32 pour
$\mathit{racine}$ et 32 pour $k$, et $n$ court jusqu'à la fin de la chaîne. Un champ de
domaine sans séparateur est donc délimité par le premier séparateur, et les champs
restants se lisent à des décalages fixes : la chaîne d'octets détermine le tuple.
\end{proof}

\noindent
Les noms d'hôte ne contiennent pas l'octet séparateur, mais c'est une propriété de
l'entrée de l'appelant et non du protocole : le portefeuille comme le vérifieur rejettent
donc un domaine qui en contiendrait un avant de hacher quoi que ce soit. Le jeton est $(k,
e, D, \pi)$.

\subsection{Vérification (hors ligne)}
\label{sec:verification}

Le site vérifie que $n$ est un nonce qu'il a émis, non expiré (TTL de 600\,s) et non
consommé ; que $D$ le désigne lui-même et que $|e - e_{\mathrm{now}}| \le 1$, où
$e_{\mathrm{now}}$ est l'époque de sa propre horloge ; recalcule $c$ et les sorties
publiques attendues $(k, \mathit{racine}, 0^d)$, où $0^d$ désigne le vecteur des $d$ zéros
que les contraintes de booléanité doivent produire, avec la racine de son
dernier lot vérifié ; puis exécute le vérifieur de preuve : pour chaque répétition,
dériver le défi depuis la transcription, redévelopper les deux graines ouvertes,
recalculer la vue entière de la première partie ouverte et l'évaluation des fils de la
seconde depuis ses sorties de multiplication engagées, et contrôler les deux engagements,
les parts de sortie des deux parties et la reconstruction à trois parts de chaque sortie
publique. En cas de succès, il marque $n$ consommé et peut associer $k$ à une session ou à
un compte. Aucun message ne quitte le site.

\subsection{Révocation et époques}

La révocation (portefeuille volé, inscription erronée, décès) est une remise de la feuille
à $E$ et un nouveau lot ; le secret révoqué ne satisfait plus $S$ contre la nouvelle
racine, et les preuves de tous les autres citoyens continuent de se vérifier une fois leur
portefeuille resynchronisé. Les époques bornent la durée de vie des pseudonymes : la clé
$k$ d'un site tourne tous les 30 jours, ce qui plafonne le profilage à long terme sur un
même site, tandis qu'au sein d'une époque le $k$ stable donne aux sites un identifiant de
visiteur récurrent sans cookies ni comptes. Les sites acceptent les jetons dont l'époque
s'écarte d'au plus un de leur époque courante (la précédente ou la suivante, conformément
au contrôle $|e - e_{\mathrm{now}}| \le 1$ de la Section~\ref{sec:verification}), tolérant
les décalages d'horloge et les effets de bord au changement d'époque.

\subsection{Seuils d'âge multiples}
\label{sec:thresholds}

Rien dans la construction n'est propre à la majorité ; la majorité n'est que le prédicat
que l'AG vérifie à l'inscription. Une juridiction dont les lois conditionnent différents
services à différents âges (13, 15, 16, 18) exploite un accumulateur par seuil : l'AG
maintient des arbres parallèles $\tree_{13}, \tree_{15}, \ldots$, un par seuil $a$, chacun avec
sa propre racine $\mathit{racine}_a$, insère l'engagement d'un citoyen dans chaque arbre
dont il satisfait le seuil, et publie un lot signé par arbre avec $a$ lié sous la
signature, de sorte qu'une racine ne peut être présentée comme certifiant un autre seuil
que celui pour lequel elle a été publiée. Le même secret de
portefeuille et le même engagement servent tous les arbres, mais les arbres se remplissent
indépendamment, si bien que le portefeuille conserve un indice de feuille par arbre ; les
insertions des citoyens qui
franchissent un seuil sont regroupées avec la cadence normale de publication des lots ; la
révocation efface la feuille dans chaque arbre. Tous les théorèmes de sécurité s'appliquent
inchangés arbre par arbre, et l'ensemble d'anonymat d'une preuve contre $\tree_a$ est
l'ensemble des citoyens inscrits satisfaisant le seuil $a$, un ensemble à l'échelle de la
population pour chaque $a$.

\paragraph{L'attaque par encadrement.} Un site dont l'obligation porte sur le seuil $a$
apprend un bit, âge $\ge a$. Rien n'oblige un site hostile à s'arrêter là. Un site qui
admet ses visiteurs à 13 ans pourrait tester chacun d'eux, en plus, contre le seuil de 18
ans : les adultes passent, un jeune de quatorze ans échoue, et le site a trié en tranches
d'âge les utilisateurs mêmes qu'il admet. Une conception multi-seuils doit fermer ce
canal, et le fermer par les mathématiques plutôt que par des conditions d'utilisation.
Trois mécanismes s'en chargent. Les deux premiers sont cryptographiques ; le troisième
rend visible, à la personne qu'elle vise, la seule demande qu'ils ne peuvent pas
éliminer.

\begin{enumerate}[leftmargin=2em]
\item \emph{Liaison du seuil.} Le seuil est soudé à la preuve elle-même. Le contexte
devient $c_a = \CtxHash(D, e, a)$ (réalisé en injectant $a$ dans la chaîne de contexte
hachée de la Section~\ref{sec:issuance}), le pseudonyme devient $k = \PRF(s, c_a)$, et la
transcription de Fiat-Shamir lie $(D, e, \mathit{racine}_a, k, n)$ ainsi que $a$. Une
preuve est donc une preuve \emph{pour un seuil} : la retester contre la racine de tout
autre seuil échoue, pour chaque visiteur, qu'il appartienne ou non à l'autre arbre
(Proposition~\ref{prop:isolation}, Section~\ref{sec:isolation}). Un retest qui échoue
pour tout le monde ne révèle rien sur personne.
\item \emph{Usage unique.} La transcription lie le nonce frais du site et le vérifieur le
consomme (Théorème~\ref{thm:replay}), si bien qu'un jeton reçu ne peut être rejoué contre
quoi que ce soit, y compris une autre racine. Le sondage ne peut donc pas recycler du
matériel que le site détient déjà ; il lui faudrait demander au portefeuille une preuve
entièrement nouvelle.
\item \emph{Transparence côté portefeuille.} Cette demande restante s'adresse au
portefeuille, et le portefeuille est l'agent du citoyen, pas celui du site. Le
portefeuille affiche quel seuil un site lui demande de prouver, retient le premier seuil
que chaque domaine a déclaré, et traite tout changement, ou toute demande de prouver un
second seuil, comme suspect : il s'interrompt avec un avertissement explicite, refuse par
défaut, et ne continue que sur confirmation délibérée de l'utilisateur (un déploiement
peut réserver ce contournement aux réglages avancés). La logique d'avertissement
s'exécute avant tout calcul dépendant de l'appartenance, elle se déclenche donc à
l'identique chez tous les portefeuilles, quel que soit l'âge du détenteur ; et un
portefeuille ne détenant aucune feuille dans l'accumulateur demandé refuse localement au
lieu d'émettre une preuve qui échouerait, si bien que le site ne peut distinguer ce cas
d'un refus. Le site ne
peut pas lire un âge dans la réaction du portefeuille, seulement dans une réponse que
l'utilisateur choisit sciemment de donner, et la demande de sondage elle-même devient une
preuve d'abus visible par l'utilisateur et consignable.
\end{enumerate}

Le seuil d'un domaine devient ainsi une partie de son identité publique. Un site reste
libre de déclarer 18 pour un service nominalement 13+, mais il refoule alors les mineurs
uniformément et visiblement, et ne reçoit toujours qu'un bit. La cascade qui semble
naturelle (tenter $\tree_{18}$, se rabattre sur $\tree_{16}$, conclure $16 \le \mathrm{\hat{a}ge}
< 18$) est sans objet pour les sites honnêtes, puisqu'un visiteur prouve directement
contre l'unique seuil déclaré du site, et constitue exactement le motif de double demande
que le mécanisme 3 signale. Comme les pseudonymes dérivent du contexte incluant le seuil,
les jetons de seuils différents sur un même domaine sont de surcroît inchaînables (le
Théorème~\ref{thm:unlink} s'applique, puisque $c_a \ne c_b$ pour $a \ne b$), défense en
profondeur si un utilisateur venait un jour à confirmer un second seuil. L'interdiction
légale de la Section~\ref{sec:deploy}, point (iii), demeure comme filet de
responsabilité, mais le mécanisme n'en dépend pas. L'implémentation de référence applique
tout cela : les lots portent le seuil sous la signature de l'AG, le haché de contexte
incorpore le seuil, les portefeuilles maintiennent un registre d'épinglages persistant et
lèvent un avertissement explicite, indépendant de l'âge, à chaque changement, et un site
configuré avec un seuil requis refuse les lots de tout autre seuil.

\section{Analyse de sécurité}
\label{sec:security}

Dans toute cette section, les fonctions de hachage sont les constructions Poseidon2 de la
Section~\ref{sec:prims}, la permutation sous-jacente est modélisée comme une permutation
aléatoire publique, SHA3/SHAKE sont modélisés comme des oracles aléatoires \cite{bellrog},
$\negl$ désigne une fonction négligeable du paramètre de sécurité, $q$ borne les requêtes
d'oracle de l'adversaire, $\tau$ est le nombre de répétitions, et $\adv$ est un
adversaire polynomial probabiliste (ou quantique, Section~\ref{sec:pq}). L'\emph{arbre
publié} désigne l'arbre binaire complet de profondeur $d$ déterminé par la suite de
feuilles publiée, toute position inoccupée portant la feuille vide $E$ et tout sous-arbre
inoccupé son empreinte mise en cache ; chaque position de chaque niveau est ainsi bien
définie.

\subsection{Infalsifiabilité (G1, R1)}

\begin{theorem}
\label{thm:unforge}
Supposons les engagements de vue liants (SHA3-256 comme oracle aléatoire), $\LeafHash$ et
$\Compress$ résistants aux collisions et à la préimage, les défis dérivés par
Fiat-Shamir dans le modèle de l'oracle aléatoire, et la racine contre laquelle le site
vérifie issue d'un lot dont il a contrôlé la signature ML-DSA-65 (de sorte que
l'``arbre publié'' soit bien défini ; sans l'EUF-CMA du schéma de signature, un lot
falsifié laisserait l'adversaire choisir la racine). Si une coalition quelconque de citoyens,
de sites et de détenteurs de l'information publique de l'AG amène un site honnête à
accepter un jeton, alors sauf avec probabilité au plus $q(2/3)^\tau + O(q^2)\,2^{-256}$ il
existe un témoin extractible $(s, \mathrm{chemin})$ tel que $\LeafHash(s)$ soit une
feuille de l'arbre publié ; et si aucun membre de la coalition n'a été inscrit avec cette
feuille, l'extraction fournit de surcroît une préimage ou une collision des fonctions de
hachage.
\end{theorem}

\begin{proof}
Fixons un jeton accepté $(k, e, D, \pi)$, soit $\mathit{racine}$ la valeur contre
laquelle le site a vérifié, et soit $c = \CtxHash(D, e)$ le contexte qu'il a recalculé.

\emph{Étape 1 (les engagements lient les vues).} Chaque répétition porte trois engagements
SHA3-256 sur (graine, vue). Produire deux ouvertures distinctes d'un même engagement exige
une collision dans l'oracle aléatoire, donc sur l'ensemble des $q$ requêtes de
$\adv$ cela survient avec probabilité au plus $q^2 2^{-256}$. Hors de cet
événement, à l'instant où une transcription est soumise pour la première fois à l'oracle
de défi, chaque répétition de cette transcription est liée à un unique triplet de vues :
deux d'entre elles entièrement déterminées par leurs graines, la troisième par sa graine
jointe à ses parts d'entrée et sorties de multiplication engagées.

\emph{Étape 2 (un énoncé faux corrompt au moins une paire par répétition).} Considérons
une répétition, son triplet de vues lié, et les trois vecteurs de parts de sortie inclus
dans la transcription. Les indices de partie parcourent $\{0,1,2\}$ et sont lus modulo 3
dans tout ce qui suit. Disons que la paire $(i, i{+}1)$ est \emph{cohérente} si le
recalcul déterministe du vérifieur pour le défi $i$ réussit : la vue de la partie $i$ se
redérive exactement de sa graine (et de ses entrées engagées le cas échéant), chaque porte
de multiplication de la partie $i$ est cohérente avec les rubans des parties $i$ et
$i{+}1$ et les sorties de porte engagées de la partie $i{+}1$, et les sorties de circuit
des deux parties égalent les parts de sortie de la transcription. Si les trois paires sont
cohérentes, les trois vecteurs de parts d'entrée définissent conjointement un témoin $w$
(la somme coordonnée par coordonnée des parts), et par récurrence sur les portes du
circuit, les trois parts de chaque fil somment à l'évaluation native de ce fil sur $w$ ;
comme le vérifieur contrôle aussi que les parts de sortie somment à $(k, \mathit{racine},
0^d)$, $w$ satisfait $S(\mathit{racine}, k, c)$. Par contraposée : si les vues liées ne
contiennent aucun témoin de $S$, au moins une des trois paires est incohérente, et le
vérifieur n'accepte la répétition que si le défi tombe sur l'une des au plus deux paires
cohérentes, probabilité au plus $2/3$.

\emph{Étape 3 (Fiat-Shamir).} Les défis d'une transcription $T$ sont la sortie de l'oracle
aléatoire sur $T$. Pour toute transcription $T$ liée comme à l'étape 1 à des vues sans
témoin, la probabilité que son vecteur de défis évite chaque paire incohérente sur les
$\tau$ répétitions est au plus $(2/3)^\tau$. Une borne d'union sur les $q$ requêtes de
l'adversaire borne sa probabilité totale de succès par $q\,(2/3)^\tau$, soit moins de
$q\,2^{-128}$ pour $\tau = 219$. L'union porte sur les transcriptions et non sur les
répétitions parce que le vérifieur dérive les $\tau$ défis du haché unique d'une seule
transcription et rejette si une seule répétition échoue : l'adversaire obtient un tirage
du vecteur de défis entier par requête, non un tirage par répétition.

\emph{Étape 4 (extraction et réduction à la sécurité des hachés).} Dans le cas restant,
les vues liées d'une transcription acceptée contiennent un témoin $(s, (b_i,
\sigma_i)_{i<d})$ de $S$. Les sorties de booléanité forcent chaque $b_i \in \{0,1\}$
exactement, donc $\MerkleFold$ effectue un authentique recalcul de racine le long du
chemin sélectionné par les bits. Comparons ce recalcul à l'arbre publié, niveau par niveau
depuis la racine : soit chaque n\oe ud produit égale le n\oe ud publié à la position
correspondante, auquel cas $\LeafHash(s)$ égale la feuille publiée à l'indice $\sum_i b_i
2^i$ ; soit il existe un niveau le plus haut où la paire d'enfants recalculée diffère des
enfants publiés alors que le parent concorde, ce qui exhibe une collision de $\Compress$.
Dans le premier cas, si la position indexée porte la feuille vide $E$, alors $(s,
\mathit{tag}_{\mathrm{leaf}}, 0, 0, 0)$ et $(0,0,0,0,\mathit{tag}_{\mathrm{empty}},0,0,0)$
sont des entrées distinctes de la permutation (elles diffèrent dans l'emplacement
d'étiquette, puisque $\mathit{tag}_{\mathrm{leaf}} \ne \mathit{tag}_{\mathrm{empty}}$, quel
que soit $s$) aux sorties tronquées égales, une collision
entre les instances à séparation de domaine ; si la position porte l'engagement d'un
citoyen honnête, $s$ en est une préimage. Chacun de ces événements contredit la sécurité
supposée des hachés, sauf avec probabilité négligeable ; un jeton accepté implique donc la
connaissance d'un secret inscrit, dans les bornes annoncées.
\end{proof}

\subsection{Anonymat et résistance à la collusion (G2, R2, R4)}

\begin{definitionx}[Jeu d'anonymat]
L'adversaire contrôle l'AG et tous les sites, mène l'inscription de citoyens de son choix
plus deux portefeuilles honnêtes $W_0, W_1$ qu'il inscrit (apprenant leurs feuilles et
indices, mais pas leurs secrets). Il interroge ensuite au plus $q$ fois un oracle de
jetons : une requête $(j, D, e, n)$ avec $j \in \{0,1\}$ renvoie le jeton que $W_j$
produirait pour $(D, e, n)$ contre la racine publiée courante. Au moment du défi, $\adv$
choisit $(D^*, e^*)$ pour lesquels il n'a fait aucune requête de jeton, pour l'un ou
l'autre portefeuille honnête, ainsi qu'un nonce $n^*$ de son choix ; le challenger tire une pièce équilibrée $\beta \in \{0,1\}$, et
$W_\beta$ produit un jeton pour $(D^*, e^*, n^*)$ que $\adv$ reçoit. $\adv$
sort une supposition $\beta'$. Le système est anonyme si $|\Pr[\beta' = \beta] - 1/2| \le
\negl$.
\end{definitionx}

\begin{theorem}
\label{thm:anon}
Dans le modèle de l'oracle aléatoire et de la permutation aléatoire, le protocole est
anonyme au sens du jeu ci-dessus, avec un avantage adversarial au plus $O(q)\,2^{-253}$.
\end{theorem}

\begin{proof}
Nous procédons par argument hybride sur le jeton de défi.

L'hybride $H_0$ est le jeu réel.

L'hybride $H_1$ remplace la preuve $\pi$ par une preuve simulée. Le simulateur tire
d'abord le vecteur de défis ; puis, pour chaque répétition, il construit les deux vues à
ouvrir exactement comme un prouveur honnête (rubans uniformes développés depuis des
graines fraîches, parts d'entrée dérivées des rubans ou explicites selon les indices de
partie), complète leurs valeurs porte par porte pour que tous les contrôles de paire et
les deux vecteurs de parts de sortie soient cohérents avec les sorties publiques, engage
honnêtement ces deux vues, engage une chaîne uniformément aléatoire à la place de la
troisième vue, et programme l'oracle aléatoire pour que la transcription résultante
corresponde aux défis pré-tirés. Les données ouvertes dans $H_1$ sont distribuées
identiquement à $H_0$ : dans les deux mondes, les deux vues ouvertes consistent en valeurs
de ruban uniformes et indépendantes et en parts dont la distribution jointe est fixée par
les mêmes contraintes de cohérence, car deux quelconques des trois parts additives de
chaque fil sont uniformes et indépendantes du témoin. Deux événements d'échec séparent les
hybrides : l'adversaire interroge l'oracle aléatoire sur la préimage de l'engagement non
ouvert, qui contient une graine fraîche de 256 bits qu'il ne voit jamais, probabilité au
plus $q\,2^{-256}$ ; ou la transcription à programmer a déjà été interrogée, ce qui exige
de deviner une des graines fraîches contenues dans les engagements, de nouveau au plus
$q\,2^{-256}$.

L'hybride $H_2$ remplace $k = \PRF(s_\beta, c^*)$, où $c^* = \CtxHash(D^*, e^*)$ et
$s_\beta$ est le secret de $W_\beta$, par une empreinte uniforme. Toutes les preuves étant
simulées après $H_1$, la vue entière que l'adversaire a de $s_\beta$ est une liste
d'empreintes : $\mathit{feuille}_\beta = \Trunc\bigl(P(s_{\beta,0}, \ldots, s_{\beta,3},
\mathit{tag}_{\mathrm{leaf}}, 0, 0, 0)\bigr)$ issue de l'inscription, les pseudonymes
$\PRF(s_\beta, c_1), \ldots, \PRF(s_\beta, c_m)$ renvoyés par ses $m \le q$ requêtes de
jetons, et le défi $k$. Notons $y = P(s_{\beta,0}, \ldots, s_{\beta,3},
\mathit{tag}_{\mathrm{prf}}, 0, 0, 0)$ l'image interne commune.

Deux observations donnent l'hybride. D'abord, $s_\beta$ est uniforme sur $\Fp^4$
(min-entropie supérieure à $255$ bits), donc l'état de feuille et l'état interne du PRF
sont deux points distincts de $P$, différant à l'emplacement d'étiquette ; $P$ étant une
permutation aléatoire publique, leurs images sont indépendantes et uniformes tant que
$\adv$ n'interroge pas $P$ ou $P^{-1}$ en un point dont la moitié de débit vaut $s_\beta$,
ce que chacune de ses $q$ requêtes fait avec probabilité au plus $2^{-255}$. Ensuite, le
jeu interdit toute requête de jeton en $(D^*, e^*)$, donc $c_i \ne c^*$ pour tout $i$ sauf
collision de $\CtxHash$ sur des entrées distinctes, probabilité au plus $q\,2^{-255}$ ;
les appels externes de la permutation sont alors évalués aux $m+1$ points distincts $y
\boxplus c_1, \ldots, y \boxplus c_m, y \boxplus c^*$, dont les images sont conjointement
uniformes et indépendantes tant que $\adv$ n'interroge pas $P$ en l'un d'eux, et $y$ est
lui-même non interrogé et uniforme. En conditionnant sur l'absence de ces événements, $k$
est uniforme et indépendant de tout le reste de ce que détient $\adv$, exactement comme
dans $H_2$. Donc $H_1$ et $H_2$ diffèrent d'au plus $O(q)\,2^{-253}$. Ce sont les $m$
pseudonymes antérieurs qui rendent cette étape nécessaire : ce sont les seuls autres
endroits où $s_\beta$ intervient, et ils sont traités plutôt qu'exclus.

Dans $H_2$ le jeton de défi est (un $k$ uniforme, $e^*$, $D^*$, un $\pi$ simulé) et
constitue un échantillon d'une distribution qui ne dépend pas de $\beta$ : l'avantage y
est donc nul. La somme des distances entre hybrides démontre le théorème.
\end{proof}

\begin{corollary}[Ce que détient une coalition en collusion totale]
La coalition de l'AG et de tous les sites possède : le registre (citoyen $\to$ feuille),
l'arbre public, et par site un ensemble de pseudonymes $k$ avec leurs métadonnées de
session. Le Théorème~\ref{thm:anon} compare deux portefeuilles ; un hybride standard sur
la population inscrite, remplaçant les pseudonymes d'un citoyen à la fois par des
empreintes uniformes, l'étend en l'énoncé que la distribution jointe de tous les
pseudonymes est calculatoirement indépendante du registre, pour un coût linéaire en le
nombre de citoyens inscrits. Il n'y a de surcroît aucun trafic à corréler :
le protocole ne définit aucun message du citoyen ou du site vers l'AG après l'inscription,
si bien que les attaques par synchronisation et par intersection dans le style de
l'analyse des réseaux de mixage n'ont aucune observation côté gouvernement à exploiter. Le
signal résiduel est celui, côté site, que possède tout service web (adresse IP, horaires),
hors périmètre par la Section~\ref{sec:model} et inchangé par ce système.
\end{corollary}

\subsection{Non-chaînabilité entre sites et entre époques (G3, R3)}

\begin{theorem}
\label{thm:unlink}
Soient $(D, e) \ne (D', e')$ deux couples domaine-époque distincts et considérons le jeu
suivant. L'adversaire contrôle l'AG et tous les sites et inscrit deux portefeuilles
honnêtes de secrets $s_0, s_1 \in \Fp^4$ ; le challenger tire $\beta \in \{0,1\}$ et rend
le jeton de $s_0$ en $(D, e)$ accompagné du jeton de $s_\beta$ en $(D', e')$, c'est-à-dire
les pseudonymes $k = \PRF(s_0, \CtxHash(D, e))$ et $k' = \PRF(s_\beta, \CtxHash(D', e'))$
avec leurs preuves. Sous le modèle du Théorème~\ref{thm:anon}, aucun adversaire ne décide
si les deux jetons proviennent d'un seul secret ($\beta = 0$) ou de deux ($\beta = 1$)
avec un avantage supérieur à $O(q)\,2^{-253}$, où $q$ borne ses requêtes à l'oracle et à
la permutation.
\end{theorem}

\begin{proof}
Posons $c = \CtxHash(D, e)$ et $c' = \CtxHash(D', e')$. Comme $\CtxHash$ est SHAKE256 avec
échantillonnage par rejet, modélisé en oracle aléatoire, ses sorties sur les deux entrées
distinctes $(D, e) \ne (D', e')$ sont indépendantes et uniformes sur $\Fp^4$, donc $c \ne
c'$ sauf avec probabilité $p^{-4} < 2^{-255}$. Il s'agit d'une collision entre deux
entrées fixées : la borne ne dépend pas du nombre de requêtes de l'adversaire.
Appliquons la séquence
hybride du Théorème~\ref{thm:anon} aux deux jetons simultanément : l'étape de simulation
est par jeton et inchangée ; l'étape de remplacement substitue aux deux pseudonymes des
valeurs uniformes, ce que justifie le même argument de permutation aléatoire, puisque les
deux appels externes de PRF évaluent $P$ en des points distincts (leurs moitiés de débit
diffèrent par $c$ contre $c'$) et sont donc conjointement uniformes et indépendants à
moins que l'adversaire n'interroge $P$ en un point dont la moitié de débit vaut $s_0$ ou
$s_1$, ce que chacune de ses $q$ requêtes fait avec probabilité au plus $2^{-255}$. Pour
$\beta = 0$ les deux appels externes partagent une même image interne et pour $\beta = 1$
ils en utilisent deux, mais dans les deux cas les points d'évaluation sont distincts :
c'est pourquoi la distinction de cas n'affleure jamais dans la vue de l'adversaire. Dans
l'hybride final,
les deux jetons sont des échantillons uniformes indépendants, une distribution identique
sous $\beta = 0$ et sous $\beta = 1$, donc l'avantage y est nul ; la somme des distances
entre hybrides donne la borne annoncée. Au sein d'un même $(D, e)$, la chaînabilité de deux visites est
intentionnelle ($k$ est déterministe) et annoncée comme la propriété de pseudonymie.
\end{proof}

\subsection{Résistance au rejeu (G4)}

\begin{theorem}
\label{thm:replay}
Soit $(k, e, D, \pi)$ un jeton produit par un portefeuille honnête et accepté par le site
$D$ avec le nonce $n$. Sauf avec probabilité au plus $(q^2 + 3\tau)\,2^{-256} +
3^{-\tau}$, le jeton ne se vérifie sur aucun autre tuple $(D', e', \mathit{racine}', k',
n') \ne (D, e, \mathit{racine}, k, n)$, et sur $(D, n)$ lui-même au plus une fois.
\end{theorem}

\begin{proof}
Écrivons la preuve comme un corps de transcription $T$ (tous les engagements et parts de
sortie) et des ouvertures $O$ pour le vecteur de défis $\chall \in \{0,1,2\}^\tau$, dérivé
comme $\chall = \RO(\mathrm{prefix} \,\|\, \mathit{ctx} \,\|\, T)$, où $\RO$ est l'oracle
aléatoire instanciant le haché de Fiat-Shamir (SHA3-256), $\mathrm{prefix}$ la chaîne fixe
de séparation de domaine du protocole, et où $\mathit{ctx} = \Enc(D, e, \mathit{racine}, k, n)$ ; $\chall_i$ désigne la $i$-ième composante de $\chall$ et les indices de
partie sont pris modulo 3. La génération honnête tire les trois graines de chaque
répétition indépendamment et uniformément, donc au sein d'une répétition les trois graines
sont distinctes deux à deux sauf avec probabilité au plus $3\,2^{-256}$, soit au plus
$3\tau\,2^{-256}$ sur les $\tau$ répétitions. Un
vérifieur pour un tuple différent dérive $\chall' = \RO(\mathrm{prefix} \,\|\, \mathit{ctx}'
\,\|\, T)$. Les deux contextes encodent des tuples distincts, donc $\mathit{ctx}' \ne \mathit{ctx}$
en tant que chaînes de bits par le Lemme~\ref{lem:enc}, et $\chall' = \chall$
seulement si l'oracle aléatoire entre en collision sur des entrées distinctes, probabilité
au plus $q^2 2^{-256}$ sur les requêtes de l'adversaire ; sinon $\chall'$ est uniforme sur
l'espace de défis à $\tau$ composantes et indépendant de $\chall$. Les ouvertures $O$
contiennent, pour chaque répétition $i$, les graines des parties $\chall_i$ et $\chall_i + 1$ et les
sorties de multiplication de la partie $\chall_i + 1$. Si $\chall'_i \ne \chall_i$ pour une répétition,
le vérifieur interprète les graines fournies comme celles des parties $\chall'_i$ et $\chall'_i + 1$
et les contrôle contre les engagements à ces indices ; chaque engagement contient la
graine fraîche de 256 bits de sa propre partie, et ces graines sont distinctes deux à
deux, donc une graine engagée à un indice
n'ouvre l'engagement d'un autre indice que par collision d'oracle aléatoire, déjà comptée.
L'acceptation exige donc $\chall' = \chall$ exactement, un événement de probabilité $3^{-\tau}$,
soit moins de $2^{-347}$ pour $\tau = 219$. La re-présentation sur $(D, n)$ lui-même est
bloquée par l'état local de nonces consommés du vérifieur, qui n'exige aucune coordination
avec quiconque ; et un nonce frais sur le même site change $\mathit{ctx}$, ce qui ramène
au cas précédent.
\end{proof}

\paragraph{Remarque.} Un portefeuille qui s'écarte de l'aléa honnête (par exemple en
donnant aux trois parties des graines identiques) peut fabriquer des jetons qui se
vérifient sous plusieurs vecteurs de défis et deviennent donc rejouables par des tiers.
Cela n'affaiblit que la liaison de ses propres jetons : cela n'accorde d'attestations de
majorité qu'avec la coopération active du détenteur du portefeuille, exactement le cas du
partage volontaire de la Section~\ref{sec:limits}.

\subsection{Isolation des seuils (déploiements multi-seuils)}
\label{sec:isolation}

Cette sous-section prouve l'affirmation du mécanisme 1 de la
Section~\ref{sec:thresholds} : un jeton n'est utilisable que contre le seuil exact pour
lequel il a été produit, et le tester contre tout autre seuil produit un verdict qui ne
porte aucune information sur le visiteur.

\begin{proposition}[Isolation des seuils]
\label{prop:isolation}
Soient $a \ne b$ des seuils d'âge de racines publiées $\mathit{racine}_a$,
$\mathit{racine}_b$, et soit $(k, e, D, \pi)$ un jeton produit par un portefeuille
honnête pour le seuil $a$ au domaine $D$ et à l'époque $e$ avec le nonce de site $n$,
c'est-à-dire une preuve de $S(\mathit{racine}_a, k, c_a)$ avec $c_a = \CtxHash(D, e, a)$,
liée par transcription comme en Section~\ref{sec:thresholds}. Pour tout nonce $n'$, un
vérifieur qui contrôle ce jeton contre le seuil $b$, c'est-à-dire contre
$(\mathit{racine}_b, c_b = \CtxHash(D, e, b), n')$, accepte avec probabilité au plus
$(q^2 + 3\tau)\,2^{-256} + 3^{-\tau}$, où $q$ borne les requêtes de l'adversaire à
l'oracle aléatoire et $\tau$ est le nombre de répétitions (Section~\ref{sec:security}).
La borne vaut uniformément sur tous les portefeuilles ; en particulier elle ne dépend pas
de la présence de la feuille du portefeuille dans $\tree_b$.
\end{proposition}

\begin{proof}
L'encodage du seuil dans le contexte de Fiat-Shamir est injectif. Les deux côtés dérivent
le contexte de leur propre $(D, a)$ en formant la chaîne $D \,\|\, \texttt{\#} \,\|\, a$
et en l'employant à la place du domaine nu, et \texttt{\#} n'apparaît dans aucun nom
d'hôte, donc l'application $(D, a) \mapsto D \,\|\, \texttt{\#} \,\|\, a$ est injective et
$a \ne b$ force des champs de domaine distincts. Ainsi la chaîne de contexte du prouveur
$\mathit{ctx}_a$, qui encode $(D \,\|\, \texttt{\#} \,\|\, a,\, e,\, \mathit{racine}_a,\,
k,\, n)$, et celle du vérifieur $b$, $\mathit{ctx}_b$, qui encode $(D \,\|\, \texttt{\#}
\,\|\, b,\, e,\, \mathit{racine}_b,\, k,\, n')$, diffèrent en tant que chaînes de bits
pour tout choix de $n'$ : elles diffèrent déjà dans le champ de domaine, que
$\mathit{racine}_a$ et $\mathit{racine}_b$ coïncident ou non. Le Lemme~\ref{lem:enc} donne alors $\mathit{ctx}_a \ne \mathit{ctx}_b$, et le
Théorème~\ref{thm:replay} s'applique mot pour mot en prenant comme tuple cible celui du
vérifieur $b$, ce qui donne la borne annoncée. Sa démonstration n'utilise que la
génération honnête de $\pi$ (graines uniformes indépendantes) et les propriétés de
l'oracle aléatoire ; elle ne fait à aucun moment référence au témoin du portefeuille ni au
contenu de $\tree_b$. Ainsi le jeton de tout portefeuille est rejeté par le vérifieur $b$
avec probabilité au moins $1 - (q^2 + 3\tau)\,2^{-256} - 3^{-\tau}$ : le verdict suit la
même distribution que le détenteur soit inscrit ou non dans $\tree_b$, et un test dont la
distribution de sortie ne dépend pas du détenteur ne transmet rien sur le détenteur. C'est
l'uniformité annoncée.
\end{proof}

\subsection{Sécurité post-quantique (G5, R5)}
\label{sec:pq}

Chaque propriété de sécurité ci-dessus se réduit à : la résistance à la préimage et la
sécurité PRF des fonctions fondées sur Poseidon2 à capacité et empreintes de 256 bits, que
Grover \cite{grover} ramène à 128 bits ; la résistance aux collisions de ces mêmes
fonctions, qui mérite le décompte séparé donné ci-dessous ; le comportement d'oracle
aléatoire de SHA3/SHAKE, où la solidité de la transformation de Fiat-Shamir face aux
adversaires quantiques découle des analyses par mesure et reprogrammation des protocoles
commit-and-open dans le modèle de l'oracle aléatoire quantique \cite{qromfs}, dont les
pertes polynomiales en requêtes sont absorbées par notre marge de solidité classique de
$2^{-128}$ ; et l'EUF-CMA de ML-DSA-65 sous Module-LWE et Module-SIS
\cite{regev,dilithium}, catégorie NIST 3. Aucune hypothèse de logarithme discret, de
factorisation ou de couplage n'apparaît nulle part. Notons le contraste avec les
déploiements pré-quantiques d'accréditations anonymes : ici, même la signature de
l'émetteur ne protège que l'authenticité de la racine ; une future cassure de ML-DSA
permettrait une falsification de racine (une attaque de solidité, atténuée par
re-signature avec un schéma successeur tel que la famille à base de hachage SPHINCS+
\cite{sphincs}) mais ne pourrait jamais désanonymiser rétroactivement qui que ce soit, car
aucune signature ne touche jamais une valeur propre à un citoyen.

\paragraph{Résistance aux collisions face au quantique.} Grover n'est pas l'outil adapté
aux collisions et nous ne revendiquons pas 128 bits de résistance quantique aux
collisions. Nos empreintes de 256 bits donnent 128 bits de résistance classique aux
collisions ; l'algorithme BHT \cite{bht} trouve des collisions en environ $2^{85}$
opérations quantiques, mais exige une mémoire quantiquement adressable comparable, et sous
toute métrique de coût qui facture cette mémoire il ne bat pas la recherche parallèle
classique \cite{bernsteincoll}. Deux observations en bornent les conséquences.
Premièrement, aucune propriété de vie privée ne repose sur la résistance aux collisions de
$\LeafHash$ ou de $\Compress$ : les Théorèmes \ref{thm:anon} et \ref{thm:unlink}
n'utilisent que le comportement de permutation aléatoire de $P$ et la min-entropie de $s$,
si bien qu'une cassure complète des collisions ne coûte aucun anonymat, et c'est l'anonymat
qui ne se répare pas après coup. Deuxièmement, les propriétés qui reposent sur la
résistance aux collisions se réparent sur place, et il vaut la peine de séparer quelle
fonction chacune sollicite. L'infalsifiabilité (Théorème~\ref{thm:unforge}) requiert la
résistance aux collisions de $\LeafHash$ et de $\Compress$, toutes deux Poseidon2 à
empreintes de 256 bits. La séparation des contextes dans le Théorème~\ref{thm:unlink}
requiert celle de $\CtxHash$, qui est une propriété de SHAKE256 tronqué dans $\Fp^4$ et
non de Poseidon2 ; les deux s'élargissent ensemble, leurs sorties faisant l'une et l'autre
quatre éléments de corps. Aucune ne descend sous le chiffre de $2^{85}$, et ce chiffre ne
vaut que sous le modèle de mémoire optimiste. Une juridiction qui juge la marge
insuffisante élargit les deux empreintes à six éléments de corps, soit environ 384 bits
(largeur d'état $t = 12$, le jeu de paramètres Poseidon2 dont nous validons déjà le
vecteur à réponse connue), pour un coût linéaire en la taille du circuit et sans changer
le protocole ni aucune des preuves ci-dessus. Notons que l'injectivité de l'\emph{encodage
du contexte} (Lemme~\ref{lem:enc}), sur laquelle reposent le Théorème~\ref{thm:replay} et
la Proposition~\ref{prop:isolation}, est une propriété syntaxique d'une disposition
d'octets : aucune avancée cryptanalytique ne l'affecte.

\subsection{Ce que le système ne résout pas}
\label{sec:limits}

\paragraph{Prêt volontaire.} Un adulte peut faire tourner le portefeuille pour un mineur,
ou lui remettre l'appareil. Aucun système d'accréditations ne l'empêche sans test
biométrique de présence continu, qui détruirait l'objectif de vie privée. Les atténuations
relèvent de la friction et de la conséquence : la rotation des époques limite le rayon
d'action d'un secret partagé ; la révocation punit le prêt commercial détecté (un secret
prêté à grande échelle est statistiquement visible pour les sites comme un pseudonyme au
comportement invraisemblable, signalable pour révocation sans identifier personne) ; la
liaison optionnelle à l'appareil (secret scellé dans un élément sécurisé) élève le coût
d'extraction. Nous rejetons consciemment l'ajout d'une étape d'autorisation
gouvernementale par usage : elle ajouterait un canal de métadonnées tandis que, étant
aveugle, elle ne pourrait de toute façon pas lier les usages aux personnes ; elle
achèterait donc un risque de surveillance sans acheter aucune résistance au partage.

\paragraph{Attaques Sybil et de solidité par l'AG.} Une AG malveillante peut inscrire des
adultes fictifs ou refuser des adultes réels. C'est une attaque contre la solidité et
l'équité, pas contre la vie privée, et elle est auditable : l'arbre est public et en ajout
seul, sa croissance doit correspondre aux statistiques démographiques publiées, et chaque
insertion peut être tenue de porter un reçu d'inclusion de journal de transparence.

\paragraph{Racines à vue divisée.} Une AG servant à une victime une racine personnalisée
est détectée par la diffusion des lots (Section~\ref{sec:rootpub}) ; un portefeuille qui
vérifie son lot contre deux miroirs indépendants réduit l'attaque à la collusion des
miroirs.

\paragraph{Coercition d'inscription et l'événement d'inscription lui-même.} L'AG sait qui
s'est inscrit. L'inscription doit donc être normalisée (groupée avec l'émission de routine
de l'eID pour tous les adultes, jamais demandée à la demande) afin que détenir une
accréditation ne porte aucune information sur l'intention de s'en servir.

\section{Implémentation}
\label{sec:impl}

L'implémentation de référence fait environ 1{,}150 lignes de Python à dépendances minimales
dans \texttt{code/python/} (\texttt{core.py} : corps, Poseidon2 en évaluation native et en
circuit, arbre de Merkle creux, prouveur et vérifieur, codec binaire de transport ;
\texttt{protocol.py} : AG, portefeuille à stockage chiffré, vérifieur de site), plus une
suite de tests et une démonstration par scénarios. Tout ce qui accompagne cet article, ses
sources comprises, se trouve sur
\url{https://github.com/Kidev/ZeroKnowledgeAgeVerification}. Points de conception
d'intérêt :

\paragraph{Circuit générique par moteur.} Le circuit est écrit une seule fois contre une
interface de moteur abstraite (const, add, sub, add-constant, mul-constant, mul) et
exécuté par trois moteurs : un moteur natif sur éléments du corps (utilisé par l'AG, le
portefeuille et le site pour la construction de l'arbre et le calcul des sorties
attendues), un moteur prouveur sur triplets de parts additives qui ajoute chaque part de
sortie de porte de multiplication aux vues par partie, et un moteur de paire qui rejoue
les deux parties défiées pour le vérifieur, recalculant la première et consommant les
sorties engagées de la seconde. La permutation native est
générique en la largeur d'état, si bien que le vecteur à réponse connue Goldilocks $t =
12$ publié par les concepteurs de Poseidon2 est contrôlé à travers cette même routine et
non à travers une seconde implémentation écrite pour le test ; la permutation du moteur de
circuit est ensuite contre-vérifiée contre elle à $t = 8$ sur des états aléatoires, et les
constantes $t = 8$ sont extraites mot pour mot de l'implémentation de référence des
concepteurs. La chaîne qui va du vecteur publié au circuit de preuve ne comporte donc
aucun maillon qu'un test pourrait satisfaire pendant que le code de production dérive.

\paragraph{Déterminisme et minimisation des vues.} Les rubans de partie sont développés
depuis des graines de 32 octets par SHAKE256 en mode compteur avec échantillonnage par
rejet (éléments du corps exactement uniformes) ; les parties 0 et 1 tirent leurs parts
d'entrée des préfixes de leurs rubans, si bien qu'une ouverture ne transmet que deux
graines, les sorties de multiplication du voisin et (quand la partie 2 est ouverte) ses
parts d'entrée. Les engagements sont SHA3-256 sur graine et vue ; le vérifieur recalcule
la vue de la partie défiée au lieu de la recevoir, si bien qu'une vue engagée corrompue
est détectée avec certitude, pas avec probabilité.

\paragraph{Format de transport et durcissement.} Les preuves se sérialisent dans un format
binaire versionné aux dispositions fixes dérivées de la profondeur et du nombre de
répétitions, chaque entier multi-octets étant en petit-boutien indépendamment de l'ordre
d'octets de la machine : le format est une spécification et non un artefact d'une
architecture. Le vérifieur est total : il contrôle les longueurs d'engagements, les plages
des parts, les longueurs de vues et la consommation exacte des rubans \emph{avant} qu'une
valeur n'atteigne une entrée de hachage, renvoie faux sur toute entrée malformée au lieu
de lever une exception, et traite le haché de transcription comme l'unique source de
vérité des défis. Le décodeur de transport borne la profondeur et le nombre de répétitions
déclarés dans l'en-tête, et les recoupe avec la longueur du blob, avant de dimensionner la
moindre allocation à partir d'eux : un décodeur qui fait confiance à une longueur fournie
par l'attaquant est une primitive d'épuisement mémoire. Les deux propriétés sont couvertes
par la suite de tests.

\paragraph{Confinement des défaillances.} Le secret du portefeuille est constitué de 4
éléments du corps uniformes générés par le CSPRNG du système, stocké uniquement sous
AES-256-GCM à clé dérivée par scrypt, et utilisé exclusivement dans $\LeafHash$ et $\PRF$
; aucun chemin de code ne le sérialise dans un message du protocole. Les chaînes de
domaine sont rejetées si elles contiennent le séparateur de contexte, ce qui acquitte
l'hypothèse du Lemme~\ref{lem:enc} au lieu de la supposer.

\paragraph{Liaison du seuil.} Dans les déploiements multi-seuils
(Section~\ref{sec:thresholds}), le seuil entre dans le lot de racine signé, dans le
contexte de preuve (encodé par la chaîne de contexte \texttt{domaine\#seuil} ; \texttt{\#}
ne peut apparaître dans un nom d'hôte, si bien que les déploiements à ensemble unique
restent compatibles octet pour octet) et dans le registre d'épinglages par domaine du
portefeuille, persistant dans le fichier de portefeuille chiffré. Le portefeuille expose
le seuil demandé pour affichage avant de prouver ; un changement de seuil lève un type
d'avertissement dédié que l'interface doit présenter, et la preuve ne se poursuit que sous
un indicateur de contournement explicite. Le contrôle d'avertissement s'exécute avant
tout calcul dépendant de l'appartenance, si bien que le comportement du portefeuille face
à un changement de seuil est indépendant de l'âge du détenteur. S'il se trouve que le
détenteur ne possède aucune feuille dans l'accumulateur demandé, le portefeuille refuse
localement au lieu d'émettre une preuve qui ne peut se vérifier : une preuve qui échoue
annoncerait ``ce détenteur est sous ce seuil'', alors qu'un refus local est
indistinguable, du point de vue du site, d'un utilisateur qui a simplement dit non. Le
site observe deux issues, jeton ou pas de jeton, au lieu de trois.

\paragraph{Portage natif.} Une implémentation Rust du corps, de la permutation, de
l'arbre, des moteurs de circuit, du prouveur et du vérifieur (\texttt{code/rust/}, environ
1\,250 lignes, partageant la transcription Fiat-Shamir et le format binaire de transport)
est arrimée à la référence par des tests à réponses connues extraits du code Python et par
l'interopérabilité du format de transport : chaque implémentation vérifie, aux paramètres
complets, des preuves produites par l'autre, si bien que les deux bases de code ne peuvent
diverger silencieusement. Ses performances mesurées sont rapportées en
Section~\ref{sec:perf}.

\section{Performances}
\label{sec:perf}

\subsection{Mesures}

Un seul c\oe ur d'un AMD Ryzen~9 7950X, CPython 3.14, Python pur (aucune extension
native), $\tau$ répétitions, arbre de profondeur $d$ :

\begin{center}
\begin{tabular}{rrrlrrr}
\toprule
profondeur $d$ & capacité & $\tau$ & solidité & preuve & vérification & taille de preuve \\
\midrule
20 & 1,05\,M  & 219 & $< 2^{-128}$ & 3,1\,s & 2,0\,s & 14,3\,Mo \\
27 & 134\,M   & 219 & $< 2^{-128}$ & 4,1\,s & 2,6\,s & 18,7\,Mo \\
32 & 4,29\,Md & 137 & $< 2^{-80}$  & 3,0\,s & 1,9\,s & 13,6\,Mo \\
32 & 4,29\,Md & 219 & $< 2^{-128}$ & 4,7\,s & 3,0\,s & 21,8\,Mo \\
\bottomrule
\end{tabular}
\end{center}

L'inscription est un haché et une insertion d'arbre en $O(d)$ (de la microseconde à la
milliseconde) ; la signature d'un lot de racine est une signature ML-DSA-65 (3\,309
octets) ; le chiffrement et le déchiffrement du portefeuille sont dominés par scrypt,
environ 100\,ms.

Le portage Rust natif du même moteur (Section~\ref{sec:impl}), produisant des preuves
identiques octet pour octet, mesure sur la même machine, sur
laquelle CPython met 4,7\,s à prouver et 3,0\,s à vérifier à $d = 32$, $\tau = 219$ :

\begin{center}
\begin{tabular}{rrrlrrr}
\toprule
profondeur $d$ & capacité & $\tau$ & solidité & preuve & vérification & taille de preuve \\
\midrule
32 & 4,29\,Md & 137 & $< 2^{-80}$  & 0,20\,s & 0,12\,s & 13,6\,Mo \\
32 & 4,29\,Md & 219 & $< 2^{-128}$ & 0,33\,s & 0,18\,s & 21,8\,Mo \\
\bottomrule
\end{tabular}
\end{center}

Soit environ quinze fois plus rapide que CPython sur la même machine, pour un portage
direct sans aucune optimisation bas niveau, et déjà sous la seconde aux paramètres
complets sur un seul c\oe ur ; les répétitions se parallélisent de surcroît entre c\oe
urs.

\subsection{Modèle de coût}

Le circuit compte $344(d+3) + 5d$ portes de multiplication ; le travail du prouveur en
vaut $3\tau$ fois, celui du vérifieur $2\tau$ fois, et la taille de preuve est dominée par
les vues voisines transmises, 8 octets par porte et par répétition. Tous les coûts sont
linéaires en $d$ et en $\tau$ et trivialement parallèles entre répétitions.

\subsection{Trajectoire de déploiement}
\label{sec:deploypath}

Les mesures ci-dessus montrent le moteur de référence franchissant déjà la barre de la
seconde dans son portage natif. Deux étapes supplémentaires, ne changeant ni le protocole
ni l'énoncé, comblent l'écart restant vers un déploiement à latence interactive. (i)
L'optimisation native du prouveur Rust inclus : Keccak vectorisé pour le développement des
rubans, évaluation du circuit en place, dispositions de parts précalculées et parallélisme
entre répétitions ramènent la génération de preuve à quelques centaines de millisecondes
et en dessous sur du matériel ordinaire. (ii) Remplacer le moteur par un STARK à base de hachage ou un système MPC-in-the-head
moderne (prétraitement de style KKW/Picnic3, ou VOLE-in-the-head) sur le circuit Poseidon2
identique réduit les preuves à la fourchette 100 à 250\,Ko et la vérification à quelques
dizaines de millisecondes, avec le même profil de confiance, sans mise en place, purement
symétrique. La vérification étant locale au site, la capacité agrégée du système est
simplement la somme des capacités des sites ; la charge de service du gouvernement est un
fichier statique (le lot et les deltas d'arbre) sur un CDN. Il n'existe aucun coût
gouvernemental par vérification, à quelque échelle que ce soit, ce qui n'est pas une
optimisation mais le point central.

\section{Considérations de déploiement}
\label{sec:deploy}

\paragraph{Lancement.} Inscrire par cohortes en marge des opérations de routine de l'eID
nationale, afin que la possession soit universelle chez les adultes et ne signale rien.
Publier l'arbre et les lots dès le premier jour avec des miroirs indépendants et un
journal d'audit en ajout seul. Diffuser le portefeuille en open source à compilation
reproductible, en magasin d'applications et en extension de navigateur ; le matériel de
clé de preuve est nul (aucune mise en place), l'intégrité de la distribution est donc la
seule préoccupation de chaîne d'approvisionnement.

\paragraph{Intégration des sites.} Un site embarque la clé de vérification des lots,
récupère les lots à intervalle régulier, émet des nonces et appelle une bibliothèque de
vérification ; toute l'intégration est sans état hormis le stockage des nonces et des
pseudonymes. Les régulateurs peuvent certifier des bibliothèques plutôt qu'auditer des
flux de données, car il n'y a pas de flux de données.

\paragraph{Cadre légal.} La législation devrait (i) autoriser l'inscription par
engagement, (ii) reconnaître un jeton vérifié comme conformité aux obligations de contrôle
d'âge, libérant la responsabilité du site, (iii) interdire aux sites d'exiger tout signal
d'identité supplémentaire quand un jeton est présenté, y compris toute preuve contre un
seuil d'âge supérieur à celui que leur obligation exige (Section~\ref{sec:thresholds}), et
(iv) rendre obligatoires les exigences de transparence de la Section~\ref{sec:rootpub}. Le point (iii) constitue un
filet de responsabilité plutôt qu'une hypothèse porteuse : le mécanisme de seuils de la
Section~\ref{sec:thresholds} rend déjà toute escalade contre-productive, de sorte que la
garantie ne dépend pas de l'interdiction.

\paragraph{Gouvernance de l'unique artefact central.} La racine est le seul objet de
confiance. Sa clé de signature devrait résider dans un HSM avec une politique publiée de
rotation et d'algorithme successeur (Section~\ref{sec:pq}), et l'historique des lots
devrait être ancré dans au moins un système hors du contrôle du gouvernement.

\section{Explication grand public}

Supposez que le gouvernement publie une longue liste publique de codes brouillés, un par
adulte. Le
brouillage est à sens unique : seul votre portefeuille, une application open source qui
tourne sur votre propre appareil, connaît le secret qui produit votre code, un secret
qu'il a généré lui-même, hors ligne, sans contacter aucun serveur. Le gouvernement sait
qu'il a
placé votre code sur la liste ; c'est l'intégralité de ce qu'il saura jamais.

Quand un site exige une preuve de majorité, le portefeuille produit une preuve à
divulgation nulle de connaissance, une technique cryptographique établie dans les années
1980 et utilisée en production aujourd'hui. Imaginez la liste du gouvernement comme une seule
chaîne fermée faite de millions de cadenas, un par adulte, et imaginez que vous portez un
masque. Vous détenez la clé d'exactement un de ces cadenas. Pour prouver votre majorité,
vous ne désignez pas votre cadenas et vous ne montrez pas votre visage : vous ouvrez
simplement la chaîne. Une chaîne ne peut s'ouvrir qu'en déverrouillant un de ses cadenas,
donc tous les témoins sont convaincus que vous détenez une clé authentique ; le masque
(ici, la couche à divulgation nulle) garantit que personne ne peut dire quel
cadenas vous avez ouvert, ni donc qui vous êtes. La démonstration est répétée quelques
centaines de fois avec de l'aléa frais, laissant à qui n'a pas de clé moins d'une chance
sur $2^{128}$ de tricher. Le site
apprend un seul fait, un adulte est présent, et reçoit un pseudonyme valable uniquement
sur ce site et pendant un mois, de quoi reconnaître un visiteur qui revient ; les
pseudonymes de sites différents sont mathématiquement sans rapport.

Pour le lecteur qui prête au gouvernement les pires intentions : après l'inscription,
votre portefeuille ne communique plus jamais avec lui, ni directement ni indirectement.
Quand
vous visitez un site, il n'existe aucun message qui puisse être intercepté, journalisé,
exigé ou divulgué, parce qu'aucun n'est envoyé. Personne n'a besoin d'être digne de
confiance pour effacer les traces de vos visites ; ces traces ne peuvent pas venir à
l'existence. Les mathématiques sont publiques, le logiciel est ouvert à l'inspection et à
la compilation reproductible, et l'unique publication du gouvernement est une liste que
chacun peut auditer contre toute édition malhonnête. Un gouvernement totalement hostile
pourrait cesser d'ajouter des adultes ou rayer des personnes de la liste, mais seulement
en public, là où chaque changement est visible. Ce qu'il ne peut pas faire, avec aucun
ordinateur y compris quantique, c'est apprendre que vous avez visité tel site. Le site, de
son côté, sait seulement qu'un adulte a présenté une preuve valide.

\section{Conclusion}

Le débat autour de la vérification d'âge en ligne présente vie privée et protection de
l'enfance comme opposées. Elles ne le sont que sous de mauvaises architectures. En
réduisant le rôle du gouvernement à la publication d'une seule liste signée d'engagements
brouillés et en repoussant toute la vérification dans des mathématiques locales entre
citoyen et site, cette conception atteint simultanément la forme forte des deux propriétés
: les sites obtiennent une attestation de majorité solide, et la correspondance entre
visites et citoyens n'existe nulle part en tant que donnée, sous collusion, sous
contrainte et sous attaque quantique. La construction n'utilise que des composants
standardisés ou académiquement établis : Poseidon2, la divulgation nulle MPC-in-the-head,
SHA-3 et ML-DSA. L'implémentation de référence accompagnant cet article démontre chaque
comportement revendiqué et chaque attaque analysée. Ce qui reste relève de l'ingénierie
(moteurs de preuve plus rapides, distribution du portefeuille) et de la législation
(imposer qu'un jeton suffise). Aucune de ces deux tâches n'exige de remettre à quiconque
la seule chose que ce système a été conçu pour ne jamais avoir à remettre.

\bibliographystyle{plain}
\bibliography{sources}

\end{document}
